\documentclass[10pt,a4paper]{amsart}
\usepackage[margin=19mm]{geometry}
\usepackage{amsmath,amssymb,mathtools}
\usepackage[scr=rsfs,cal=euler]{mathalfa}
\usepackage{xfrac}
\usepackage[unicode,hidelinks,bookmarks=false]{hyperref}
\newcommand{\ie}{i.e.\ }
\newcommand{\cf}{cf.\ }
\newcommand{\resp}{respectively\ }
\newcommand{\Subsection}{\subsection}
\providecommand{\nobreakdash}{\nobreak\hbox{-}}
\setlength{\emergencystretch}{2em}
\multlinegap0pt
\usepackage{amssymb}
\usepackage[shortlabels]{enumitem}
\usepackage{mathtools}
\newtheorem{lemma}{Lemma}
\newcommand{\eucl}{\operatorname{eucl}}
\title{Viscosity Solutions in Martinet Spaces}
\author{Thomas Bieske, Frederic Bowen}
\date{Source: arXiv:2601.19864v1 (CC BY 4.0)}
\begin{document}
\maketitle

\noindent\emph{Source and licence.} Viscosity Solutions in Martinet Spaces. Version arXiv:2601.19864v1 (CC BY 4.0), distributed by arXiv under the Creative Commons Attribution 4.0 International licence, http://creativecommons.org/licenses/by/4.0/, as recorded in the arXiv licence metadata for this exact version and shown on the abstract page https://arxiv.org/abs/2601.19864v1. Full licence notice: \texttt{licenses/arxiv-2601.19864-CC-BY-4.0.txt}.\par\smallskip

\noindent\textit{Editorial scope.} Counted unit: the article's lemma 5.2 (source0.tex lines 237--241). The article's complete original proof of it is delivered with it (source0.tex lines 243--312, one \texttt{proof} environment of 70 lines). No mathematics is written, repaired or re-derived: the statement and the proof are the article's own lines, in the article's own notation, and no retained line is substituted or re-flowed.  Retained uncounted as the source-local context the proof needs: lines 162--187.  Nothing was omitted from any retained span, so the package carries no omission record.  No retained line contains a citation, so the wrapper carries no bibliography and no bibliography entry.  No retained line contains a figure environment, an image-inclusion command or drawing code, so no image is delivered and the package declares no asset dependency.  Editorial formatting only, in addition: an AMS \texttt{amsart} preamble that restates the article's own declarations and macros used by the retained lines, namely the environments \texttt{lemma} on separate counters, together with the standard packages those lines need.  The retained environments print in this document as Lemma 1 in order of appearance, whereas the article prints the same items with the numbering of its own full document.  The complete native source of the article as distributed in the arXiv e-print for this exact version is bundled unchanged as \texttt{sources/193492/source0.tex}.
\begin{lemma}[Twisting Lemma]\label{twlemma}
Let $O\subset \mathbb{M}$ be open, let $u: O\to\mathbb{R}$, and let $p_0=(x_1^0,x_2^0,x_3^0) \in O$. Suppose that $(\eta, X)\in J^{2,+}_{\textmd{eucl}} u(p_0)$ with $\eta=(\eta_1,\eta_2,\eta_3)$ and $X=\{X_{ij}\}$ a $3\times 3$ symmetric matrix. Then 
\begin{equation}\label{twisty}
(B(p_0)\cdot \eta, A(p_0)\cdot X\cdot A^T(p_0)+\mathcal{T}(\eta,p_0))\in J^{2,+} u(p_0)
\end{equation}
where $A(p_0)$ is the $2\times3$ matrix of Martinet vector coefficients defined by
\[A(p_0)=\begin{pmatrix}1&0&0\\0&1&f(x_1^0)\end{pmatrix}_,\]
the $3\times 3$ matrix $B(p_0)$ is given by 
\[B(p_0)=A(p_0)\oplus(0,0,f'(x_1^0)=\begin{pmatrix}1&0&0\\0&1&f(x_1^0)\\0&0&f'(x_1^0) \end{pmatrix}_,\]
and the $2\times2$ twisting matrix $\mathcal{T}(\eta,p_0)$ is given by
\[\mathcal{T}(\eta,p_0)=\frac{\eta_3}{2}\begin{pmatrix}0&f'(x_1^0)\\
f'(x_1^0)&0\end{pmatrix}_.\]
More specifically,
\[B(p_0)\cdot\eta=\begin{pmatrix}
\eta_1\\
\eta_2+f(x_1^0)\eta_3\\
f'(x_1^0)\eta_3
\end{pmatrix}\]
and 
\begin{eqnarray*}
\lefteqn{A(p_0)\cdot X\cdot A^T(p_0)+\mathcal{T}(\eta,p_0)=} & & \\
& & \begin{pmatrix}x_{11}&x_{12}+f(x_1^0)x_{13}+\frac{1}{2}f'(x_1^0)\eta_3\\
\mbox{} & \mbox{} \\
x_{12}+f(x_1^0)x_{13}+\frac{1}{2}f'(x_1^0)\eta_3&x_{22}+2f(x_1^0)x_{23}+f(x_1^0)^2x_{33}\end{pmatrix}_.
\end{eqnarray*}
\end{lemma}

\begin{lemma}\label{5.2}
The vectors $\Upsilon_{p_\tau}$ and $\Upsilon_{q_\tau}$ satisfy
\begin{equation}\|\Upsilon_{p_\tau}\|^2-\|\Upsilon_{q_\tau}\|^2=O(\varphi(p_\tau,q_\tau)^2).\end{equation}
In addition, with the usual ordering, the matrix $\mathcal{X}^\tau$ is smaller than the matrix $\mathcal{Y}^\tau$ with an error term. In particular we have $\mathcal{X}^\tau\leq\mathcal{Y}^\tau+\mathcal{R}^\tau$, where $\mathcal{R}^\tau\rightarrow0$ as $\tau\rightarrow\infty$
\end{lemma}

\begin{proof}
A straightforward computation shows 
\begin{eqnarray*}
\|\Upsilon_{p_\tau}\|^2-\|\Upsilon_{q_\tau}\|^2 & = & 2(x_2-y_2)^3(x_3-y_3)^3(f(x_1)-f(y_1)) \\ 
& & \mbox{}+(x_3-y_3)^6(f(x_1)^2-f(y_1)^2)+(f'(x_1)-f'(y_1))(x_3-y_3)^6
\end{eqnarray*}
Now, \begin{equation*}
(x_2-y_2) \sim  \varphi^\frac{1}{4},\
(x_3-y_3)  \sim  \varphi^\frac{1}{4},\ \textmd{and\ }
f(x_1)-f(y_1)  \sim  (x_1-y_1)\sim\varphi^\frac{1}{2}.
\end{equation*}
We then have 
\[\|\Upsilon_{p_\tau}\|^2-\|\Upsilon_{q_\tau}\|^2\sim (\varphi^\frac{8}{4}+\varphi^\frac{8}{4}+\varphi^\frac{8}{4}).\]
The vector result follows. 

Next, we focus on the matrix difference estimate.  Using the Twisting Lemma (Lemma \ref{twlemma}), we have (writing $A_p$ for $A(p_\tau)$, $\mathcal{T}_p$ for $\mathcal{T}(\Upsilon_{p_\tau},p_\tau))$, etc),

\begin{eqnarray*}
\langle \mathcal{X}^\tau\varepsilon,\varepsilon\rangle-\langle \mathcal{Y}^\tau\kappa,\kappa\rangle
&=& \langle(A_pX^\tau A^T_p+\mathcal{T}_p)\epsilon,\epsilon\rangle-\langle(A_qY^\tau A^T_q+\mathcal{T}_q)\kappa,\kappa\rangle\\
&=&\langle X^\tau A^T_p\epsilon,A^T_p\epsilon\rangle-\langle Y^\tau A^T_q\kappa,A^T_q\kappa\rangle+\langle \mathcal{T}_p\epsilon,\epsilon\rangle-\langle \mathcal{T}_q\kappa,\kappa\rangle\\
& &\leq\tau\left\langle C
\begin{pmatrix} A^T_p\epsilon\\A^T_q\kappa \end{pmatrix},
\begin{pmatrix} A^T_p\epsilon\\A^T_q\kappa \end{pmatrix}\right\rangle +\langle \mathcal{T}_p\epsilon,\epsilon\rangle-\langle \mathcal{T}_q\kappa,\kappa\rangle \\
 & \stackrel{\textmd{def}}{=} & \tau\mathcal{V}+\mathcal{W}. 
\end{eqnarray*}
Here the matrix $C$ is given by 
\begin{eqnarray*}
\lefteqn{C=(D^2\varphi)_{\eucl}=}& & \\
& \begin{pmatrix}
1& 0 & 0 &-1 & 0 & 0 \\
0& 3(x_2-y_2)^2 & 0 & 0 & -3(x_2-y_2)^2 & 0 \\
0& 0 & 3(x_3-y_3)^2 & 0 & 0 & 3(x_3-y_3)^2 \\
-1& 0 & 0 & 1 & 0 & 0 \\
0 & -3(x_2-y_2)^2 & 0 & 0& 3(x_2-y_2)^2 & 0 \\
 0 & 0 & 3(x_3-y_3)^2 & 0& 0 & 3(x_3-y_3)^2 
\end{pmatrix}_. 
\end{eqnarray*}
Let us first simplify the inner product given by the term $\tau\mathcal{V}$. 
Unravelling the matrix multiplication, we have 
\begin{equation*}
C \begin{pmatrix} A^T_p\epsilon\\A^T_q\kappa \end{pmatrix} = 
\begin{pmatrix}
\epsilon_1-\kappa_1\\
 3(x_2^\tau-y_2^\tau)^2 (\epsilon_2-\kappa_2)\\
 3(x_3^\tau-y_3^\tau)^2(f(x_1^\tau)\epsilon_2-f(y_1^\tau)\kappa_2)\\
 -\epsilon_1+\kappa_1\\
 3(x_2^\tau-y_2^\tau)^2(-\epsilon_2+\kappa_2)\\ 
 3(x_3^\tau-y_3^\tau)^2(-f(x_1^\tau)\epsilon_2+f(y_1^\tau)\kappa_2)
 \end{pmatrix}
\end{equation*}
so that $\tau \mathcal{V}$ is given by 
\begin{equation*}
\tau\mathcal{V} = \tau\bigg((\epsilon_1-\kappa_1)^2+ 3(x_2^\tau-y_2^\tau)^2(\epsilon_2-\kappa_2)^2+  3(x_3^\tau-y_3^\tau)^2(f(x_1^\tau)\epsilon_2-f(y_1^\tau)\kappa_2)^2\bigg)
\end{equation*}
We note that when $\epsilon=\kappa$, we have 
\begin{eqnarray}\label{Vlabel}
\tau \mathcal{V} & \sim & \tau\bigg((x_3^\tau-y_3^\tau)^2(f(x_1^\tau)-f(y_1^\tau))^2\bigg) \sim \tau \varphi^{\frac{1}{2}} \times \varphi
\end{eqnarray}
We now focus on the $\mathcal{W}$ term. We have 
\begin{equation*}
\mathcal{W}=\langle \mathcal{T}_p\epsilon,\epsilon\rangle-\langle \mathcal{T}_q\kappa,\kappa\rangle =
\tau (x_3^\tau-y_3^\tau)^3\bigg(f'(x_1^\tau)\epsilon_1\epsilon_2-f'(y_1^\tau)\kappa_1\kappa_2\bigg).
\end{equation*}
When $\epsilon=\kappa$, we then have 
\begin{equation}\label{Wlabel}
\mathcal{W}\sim \tau\varphi^{\frac{3}{4}}\varphi^{\frac{1}{2}}
\end{equation}
The matrix relation then follows from Equations \eqref{Vlabel} and \eqref{Wlabel}. 
\end{proof}

\end{document}
